\documentclass[11pt,a4paper]{article}
\usepackage[utf8]{inputenc}
\usepackage[T1]{fontenc}
\usepackage[spanish,es-minimal,es-noshorthands]{babel}
\usepackage{lmodern}
\usepackage[sfdefault,lining]{FiraSans}
\usepackage[a4paper,top=2.3cm,bottom=2.2cm,left=1.8cm,right=1.8cm,headheight=1.3cm,headsep=0.4cm,footskip=1cm]{geometry}
\usepackage{amsmath,amssymb}
\usepackage{icomma}
\usepackage{graphicx}
\usepackage[table]{xcolor}
\usepackage{tikz}
\usepackage{fontawesome5}
\usepackage{ccicons}
\usepackage{fancyhdr}
\usepackage{lastpage}
\usepackage{tasks}
\usepackage{needspace}
\usepackage[most]{tcolorbox}
\usepackage{array}
\usepackage[hidelinks]{hyperref}
\AtBeginDocument{\hypersetup{pdftitle={Ficha de ejercicios de la unidad 3: Números enteros. 1º ESO},pdfsubject={Matemáticas 1º ESO},pdfkeywords={números enteros, recta numérica, valor absoluto, opuesto, suma y resta, regla de los signos, operaciones combinadas, problemas}}}

\definecolor{navy}{HTML}{1B2748}
\definecolor{accent}{HTML}{E8603A}
\definecolor{teal}{HTML}{1F9E90}
\definecolor{soft}{HTML}{EEF1F8}
\definecolor{muted}{HTML}{5D6685}

\setlength{\parindent}{0pt}
\setlength{\parskip}{3pt}
\settasks{label=\alph*), label-format=\bfseries\color{navy}, label-width=1.3em, item-indent=1.8em, column-sep=0.6em, after-item-skip=2pt, before-skip=3pt, after-skip=4pt}

% ---------- capçalera i peu ----------
\newcommand{\llicencia}{\href{https://creativecommons.org/licenses/by-nc-sa/4.0/deed.es}{\ccbyncsa}\ \ Licencia Creative Commons Atribución-NoComercial-CompartirIgual 4.0 Internacional (CC BY-NC-SA 4.0)}
\pagestyle{fancy}
\fancyhf{}
\fancyhead[L]{\raisebox{-0.35\height}{\includegraphics[height=0.95cm]{logo-gva.png}}}
\fancyhead[C]{\small\color{muted}Ficha de ejercicios · Unidad 3: Números enteros · 1\textsuperscript{o} ESO}
\fancyhead[R]{\raisebox{-0.35\height}{\includegraphics[height=0.95cm]{logo-vermellar.png}}}
\fancyfoot[L]{\scriptsize\color{muted}\llicencia}
\fancyfoot[R]{\small\color{muted}\thepage}
\renewcommand{\headrulewidth}{0.4pt}
\renewcommand{\headrule}{\color{muted!40}\hrule height \headrulewidth}
\fancypagestyle{primera}{\fancyhf{}\renewcommand{\headrulewidth}{0pt}
  \fancyfoot[L]{\scriptsize\color{muted}\llicencia}
  \fancyfoot[R]{\small\color{muted}\thepage}}

% ---------- nivells i exercicis ----------
\newcommand{\estrelles}[1]{{\color{accent}\ifcase#1\or\faStar\faStar[regular]\faStar[regular]\or\faStar\faStar\faStar[regular]\or\faStar\faStar\faStar\fi}}
\newcounter{exer}
\newcommand{\ex}[1]{\par\needspace{4\baselineskip}\medskip\stepcounter{exer}\noindent{\color{navy}\bfseries\theexer.}\ \ {\small\estrelles{#1}}\ \ }

\newcommand{\seccio}[2]{\par\needspace{10\baselineskip}\bigskip
  \begin{tcolorbox}[colback=navy,colframe=navy,coltext=white,arc=2mm,boxrule=0pt,left=3mm,right=3mm,top=1.2mm,bottom=1.2mm]
  {\large\bfseries #1\quad #2}
  \end{tcolorbox}\vspace{-2pt}}

\newcolumntype{C}[1]{>{\centering\arraybackslash}p{#1}}

% quadre de recordatori (només a la fitxa d'exercicis)
\newcommand{\recorda}[1]{\par\smallskip\begin{tcolorbox}[colback=accent!7,colframe=accent!55,arc=2mm,boxrule=0.5pt,left=2.5mm,right=2.5mm,top=1mm,bottom=1mm]\small{\color{accent}\faLightbulb[regular]}\ \textbf{Recuerda:}\ #1\end{tcolorbox}}
% recta numèrica entera: \rectabuida{min}{max} i \recta{min}{max}{punts}
\newcommand{\rectabase}[2]{\draw[latex-latex,muted] (#1-0.8,0)--(#2+0.8,0);
  \foreach \i in {#1,...,#2}{\draw[muted] (\i,0.12)--(\i,-0.12) node[below=-1pt,font=\scriptsize,text=muted]{$\i$};}}
\newcommand{\rectabuida}[2]{\begin{tikzpicture}[x=0.78cm,baseline=0]\rectabase{#1}{#2}\end{tikzpicture}}
\newcommand{\recta}[3]{\begin{tikzpicture}[x=0.78cm,baseline=0]\rectabase{#1}{#2}
  \foreach \p in {#3}{\fill[accent] (\p,0) circle (0.11) node[above=1pt,font=\small\bfseries,text=navy]{$\p$};}\end{tikzpicture}}
\newcommand{\res}[1]{#1}

% =====================================================================
% Capçalera FentMates de les fitxes: un color per a cada curs.
% Els colors són els de la pàgina inicial (LEVELS a assets/js/materials.js);
% el logo ve de _eines/logo/pdf/ (python3 _eines/logo/genera.py).
%
% Ús, al final del preàmbul de la unitat:
%   % =====================================================================
% Capçalera FentMates de les fitxes: un color per a cada curs.
% Els colors són els de la pàgina inicial (LEVELS a assets/js/materials.js);
% el logo ve de _eines/logo/pdf/ (python3 _eines/logo/genera.py).
%
% Ús, al final del preàmbul de la unitat:
%   % =====================================================================
% Capçalera FentMates de les fitxes: un color per a cada curs.
% Els colors són els de la pàgina inicial (LEVELS a assets/js/materials.js);
% el logo ve de _eines/logo/pdf/ (python3 _eines/logo/genera.py).
%
% Ús, al final del preàmbul de la unitat:
%   \input{../comu/fentmates.tex}
%   \fmcurs{1eso}
% i, al començament del document:
%   \capcalera{Matemàtiques}{1r ESO}{Fitxa d'exercicis · Unitat 2}{Divisibilitat}
% Opcional abans de \input: \def\fmneed{14} (espai mínim abans de cada secció).
% =====================================================================
\makeatletter
\def\fm@logodir{../../logo/pdf}
\providecommand{\fmneed}{10}
\def\fm@cursnom{}

% clau del curs, color i color del text damunt del color
\newcommand{\fm@curs}[3]{\@namedef{fm@color@#1}{#2}\@namedef{fm@text@#1}{#3}}
\fm@curs{1eso}{E6B84A}{navy}
\fm@curs{2eso}{3B82F6}{white}
\fm@curs{3eso}{FF7A50}{white}
\fm@curs{4eso-a}{A78BFA}{white}
\fm@curs{4eso-b}{6D28D9}{white}
\fm@curs{1batx-mi}{9B1C1C}{white}
\fm@curs{1batx-ccss}{EF6B6B}{white}
\fm@curs{2batx-mii}{1E6B47}{white}
\fm@curs{2batx-ccss}{2BB3A3}{white}

\newcommand{\fmcurs}[1]{%
  \@ifundefined{fm@color@#1}{\PackageError{fentmates}{Curs desconegut: #1}{}}{}%
  \edef\fm@tmp{\noexpand\definecolor{curs}{HTML}{\@nameuse{fm@color@#1}}%
               \noexpand\colorlet{curstext}{\@nameuse{fm@text@#1}}}\fm@tmp
  \colorlet{cursfosc}{curs!72!black}%
  \def\fm@logo{\fm@logodir/vertical-#1.pdf}%
  \def\fm@marca{\fm@logodir/marca-#1.pdf}}

% ---------- capçalera de la primera pàgina ----------
% #1 matèria (petit), #2 curs, #3 tipus de document i unitat, #4 títol
\newcommand{\capcalera}[4]{%
  \gdef\fm@cursnom{#2}%
  \par\noindent
  \begin{tcolorbox}[colback=curs,colframe=curs,boxrule=0pt,arc=2.2mm,width=3.5cm,height=2.15cm,
      valign=center,left=3mm,right=2mm,top=1mm,bottom=1mm,nobeforeafter,box align=center]
    \color{curstext}{\footnotesize #1}\par\vspace{1pt}{\fontsize{19}{21}\selectfont\bfseries #2\par}%
  \end{tcolorbox}\hspace{5mm}%
  \begin{minipage}[c]{\dimexpr\linewidth-3.5cm-5mm-2.9cm\relax}
    {\color{cursfosc}\small\bfseries\MakeUppercase{#3}}\par\vspace{2pt}
    {\color{navy}\fontsize{20}{23}\selectfont\bfseries #4\par}
  \end{minipage}\hfill
  \begin{minipage}[c]{2.6cm}\raggedleft\includegraphics[width=2.6cm]{\fm@logo}\end{minipage}\par
  \vspace{2pt}}

% ---------- barres de secció del color del curs ----------
\renewcommand{\seccio}[2]{\par\needspace{\fmneed\baselineskip}\bigskip
  \begin{tcolorbox}[colback=curs,colframe=curs,coltext=curstext,arc=2mm,boxrule=0pt,left=3mm,right=3mm,top=1.2mm,bottom=1.2mm]
  {\large\bfseries #1\quad #2}
  \end{tcolorbox}\vspace{-2pt}}

% ---------- estrelles de dificultat del color del curs (a Batxillerat no se'n posen) ----------
\@ifundefined{estrelles}{}{%
  \renewcommand{\estrelles}[1]{{\color{curs}\ifcase#1\or\faStar\faStar[regular]\faStar[regular]\or\faStar\faStar\faStar[regular]\or\faStar\faStar\faStar\fi}}}

% ---------- capçalera i peu de les pàgines ----------
% primera pàgina: logos institucionals al peu
\fancypagestyle{primera}{\fancyhf{}\renewcommand{\headrulewidth}{0pt}
  \fancyfoot[L]{\raisebox{-0.32\height}{\includegraphics[height=0.78cm]{logo-gva.png}}\hspace{5mm}%
    \raisebox{-0.32\height}{\includegraphics[height=0.78cm]{logo-vermellar.png}}}
  \fancyfoot[R]{\scriptsize\color{muted}\href{https://creativecommons.org/licenses/by-nc-sa/4.0/deed.ca}{\ccbyncsa}\ \ CC BY-NC-SA 4.0 · fentmates.org\qquad\small\thepage}}
% la resta: el símbol a l'esquerra i el curs a la dreta, amb el títol al centre com abans
\fancyhead[L]{\raisebox{-0.3\height}{\includegraphics[height=0.62cm]{\fm@marca}}}
\fancyhead[R]{\small\bfseries\color{cursfosc}\fm@cursnom}
\renewcommand{\headrule}{\color{curs!45}\hrule height \headrulewidth}
\makeatother

%   \fmcurs{1eso}
% i, al començament del document:
%   \capcalera{Matemàtiques}{1r ESO}{Fitxa d'exercicis · Unitat 2}{Divisibilitat}
% Opcional abans de \input: \def\fmneed{14} (espai mínim abans de cada secció).
% =====================================================================
\makeatletter
\def\fm@logodir{../../logo/pdf}
\providecommand{\fmneed}{10}
\def\fm@cursnom{}

% clau del curs, color i color del text damunt del color
\newcommand{\fm@curs}[3]{\@namedef{fm@color@#1}{#2}\@namedef{fm@text@#1}{#3}}
\fm@curs{1eso}{E6B84A}{navy}
\fm@curs{2eso}{3B82F6}{white}
\fm@curs{3eso}{FF7A50}{white}
\fm@curs{4eso-a}{A78BFA}{white}
\fm@curs{4eso-b}{6D28D9}{white}
\fm@curs{1batx-mi}{9B1C1C}{white}
\fm@curs{1batx-ccss}{EF6B6B}{white}
\fm@curs{2batx-mii}{1E6B47}{white}
\fm@curs{2batx-ccss}{2BB3A3}{white}

\newcommand{\fmcurs}[1]{%
  \@ifundefined{fm@color@#1}{\PackageError{fentmates}{Curs desconegut: #1}{}}{}%
  \edef\fm@tmp{\noexpand\definecolor{curs}{HTML}{\@nameuse{fm@color@#1}}%
               \noexpand\colorlet{curstext}{\@nameuse{fm@text@#1}}}\fm@tmp
  \colorlet{cursfosc}{curs!72!black}%
  \def\fm@logo{\fm@logodir/vertical-#1.pdf}%
  \def\fm@marca{\fm@logodir/marca-#1.pdf}}

% ---------- capçalera de la primera pàgina ----------
% #1 matèria (petit), #2 curs, #3 tipus de document i unitat, #4 títol
\newcommand{\capcalera}[4]{%
  \gdef\fm@cursnom{#2}%
  \par\noindent
  \begin{tcolorbox}[colback=curs,colframe=curs,boxrule=0pt,arc=2.2mm,width=3.5cm,height=2.15cm,
      valign=center,left=3mm,right=2mm,top=1mm,bottom=1mm,nobeforeafter,box align=center]
    \color{curstext}{\footnotesize #1}\par\vspace{1pt}{\fontsize{19}{21}\selectfont\bfseries #2\par}%
  \end{tcolorbox}\hspace{5mm}%
  \begin{minipage}[c]{\dimexpr\linewidth-3.5cm-5mm-2.9cm\relax}
    {\color{cursfosc}\small\bfseries\MakeUppercase{#3}}\par\vspace{2pt}
    {\color{navy}\fontsize{20}{23}\selectfont\bfseries #4\par}
  \end{minipage}\hfill
  \begin{minipage}[c]{2.6cm}\raggedleft\includegraphics[width=2.6cm]{\fm@logo}\end{minipage}\par
  \vspace{2pt}}

% ---------- barres de secció del color del curs ----------
\renewcommand{\seccio}[2]{\par\needspace{\fmneed\baselineskip}\bigskip
  \begin{tcolorbox}[colback=curs,colframe=curs,coltext=curstext,arc=2mm,boxrule=0pt,left=3mm,right=3mm,top=1.2mm,bottom=1.2mm]
  {\large\bfseries #1\quad #2}
  \end{tcolorbox}\vspace{-2pt}}

% ---------- estrelles de dificultat del color del curs (a Batxillerat no se'n posen) ----------
\@ifundefined{estrelles}{}{%
  \renewcommand{\estrelles}[1]{{\color{curs}\ifcase#1\or\faStar\faStar[regular]\faStar[regular]\or\faStar\faStar\faStar[regular]\or\faStar\faStar\faStar\fi}}}

% ---------- capçalera i peu de les pàgines ----------
% primera pàgina: logos institucionals al peu
\fancypagestyle{primera}{\fancyhf{}\renewcommand{\headrulewidth}{0pt}
  \fancyfoot[L]{\raisebox{-0.32\height}{\includegraphics[height=0.78cm]{logo-gva.png}}\hspace{5mm}%
    \raisebox{-0.32\height}{\includegraphics[height=0.78cm]{logo-vermellar.png}}}
  \fancyfoot[R]{\scriptsize\color{muted}\href{https://creativecommons.org/licenses/by-nc-sa/4.0/deed.ca}{\ccbyncsa}\ \ CC BY-NC-SA 4.0 · fentmates.org\qquad\small\thepage}}
% la resta: el símbol a l'esquerra i el curs a la dreta, amb el títol al centre com abans
\fancyhead[L]{\raisebox{-0.3\height}{\includegraphics[height=0.62cm]{\fm@marca}}}
\fancyhead[R]{\small\bfseries\color{cursfosc}\fm@cursnom}
\renewcommand{\headrule}{\color{curs!45}\hrule height \headrulewidth}
\makeatother

%   \fmcurs{1eso}
% i, al començament del document:
%   \capcalera{Matemàtiques}{1r ESO}{Fitxa d'exercicis · Unitat 2}{Divisibilitat}
% Opcional abans de \input: \def\fmneed{14} (espai mínim abans de cada secció).
% =====================================================================
\makeatletter
\def\fm@logodir{../../logo/pdf}
\providecommand{\fmneed}{10}
\def\fm@cursnom{}

% clau del curs, color i color del text damunt del color
\newcommand{\fm@curs}[3]{\@namedef{fm@color@#1}{#2}\@namedef{fm@text@#1}{#3}}
\fm@curs{1eso}{E6B84A}{navy}
\fm@curs{2eso}{3B82F6}{white}
\fm@curs{3eso}{FF7A50}{white}
\fm@curs{4eso-a}{A78BFA}{white}
\fm@curs{4eso-b}{6D28D9}{white}
\fm@curs{1batx-mi}{9B1C1C}{white}
\fm@curs{1batx-ccss}{EF6B6B}{white}
\fm@curs{2batx-mii}{1E6B47}{white}
\fm@curs{2batx-ccss}{2BB3A3}{white}

\newcommand{\fmcurs}[1]{%
  \@ifundefined{fm@color@#1}{\PackageError{fentmates}{Curs desconegut: #1}{}}{}%
  \edef\fm@tmp{\noexpand\definecolor{curs}{HTML}{\@nameuse{fm@color@#1}}%
               \noexpand\colorlet{curstext}{\@nameuse{fm@text@#1}}}\fm@tmp
  \colorlet{cursfosc}{curs!72!black}%
  \def\fm@logo{\fm@logodir/vertical-#1.pdf}%
  \def\fm@marca{\fm@logodir/marca-#1.pdf}}

% ---------- capçalera de la primera pàgina ----------
% #1 matèria (petit), #2 curs, #3 tipus de document i unitat, #4 títol
\newcommand{\capcalera}[4]{%
  \gdef\fm@cursnom{#2}%
  \par\noindent
  \begin{tcolorbox}[colback=curs,colframe=curs,boxrule=0pt,arc=2.2mm,width=3.5cm,height=2.15cm,
      valign=center,left=3mm,right=2mm,top=1mm,bottom=1mm,nobeforeafter,box align=center]
    \color{curstext}{\footnotesize #1}\par\vspace{1pt}{\fontsize{19}{21}\selectfont\bfseries #2\par}%
  \end{tcolorbox}\hspace{5mm}%
  \begin{minipage}[c]{\dimexpr\linewidth-3.5cm-5mm-2.9cm\relax}
    {\color{cursfosc}\small\bfseries\MakeUppercase{#3}}\par\vspace{2pt}
    {\color{navy}\fontsize{20}{23}\selectfont\bfseries #4\par}
  \end{minipage}\hfill
  \begin{minipage}[c]{2.6cm}\raggedleft\includegraphics[width=2.6cm]{\fm@logo}\end{minipage}\par
  \vspace{2pt}}

% ---------- barres de secció del color del curs ----------
\renewcommand{\seccio}[2]{\par\needspace{\fmneed\baselineskip}\bigskip
  \begin{tcolorbox}[colback=curs,colframe=curs,coltext=curstext,arc=2mm,boxrule=0pt,left=3mm,right=3mm,top=1.2mm,bottom=1.2mm]
  {\large\bfseries #1\quad #2}
  \end{tcolorbox}\vspace{-2pt}}

% ---------- estrelles de dificultat del color del curs (a Batxillerat no se'n posen) ----------
\@ifundefined{estrelles}{}{%
  \renewcommand{\estrelles}[1]{{\color{curs}\ifcase#1\or\faStar\faStar[regular]\faStar[regular]\or\faStar\faStar\faStar[regular]\or\faStar\faStar\faStar\fi}}}

% ---------- capçalera i peu de les pàgines ----------
% primera pàgina: logos institucionals al peu
\fancypagestyle{primera}{\fancyhf{}\renewcommand{\headrulewidth}{0pt}
  \fancyfoot[L]{\raisebox{-0.32\height}{\includegraphics[height=0.78cm]{logo-gva.png}}\hspace{5mm}%
    \raisebox{-0.32\height}{\includegraphics[height=0.78cm]{logo-vermellar.png}}}
  \fancyfoot[R]{\scriptsize\color{muted}\href{https://creativecommons.org/licenses/by-nc-sa/4.0/deed.ca}{\ccbyncsa}\ \ CC BY-NC-SA 4.0 · fentmates.org\qquad\small\thepage}}
% la resta: el símbol a l'esquerra i el curs a la dreta, amb el títol al centre com abans
\fancyhead[L]{\raisebox{-0.3\height}{\includegraphics[height=0.62cm]{\fm@marca}}}
\fancyhead[R]{\small\bfseries\color{cursfosc}\fm@cursnom}
\renewcommand{\headrule}{\color{curs!45}\hrule height \headrulewidth}
\makeatother

\fmcurs{1eso}

\begin{document}
\thispagestyle{primera}

\capcalera{Matemáticas}{1\textsuperscript{o} ESO}{Ficha de ejercicios · Unidad 3}{Números enteros}
\vspace{10pt}
Nombre y apellidos: \rule{6.6cm}{0.4pt}\hfill Grupo: \rule{1.6cm}{0.4pt}\hfill Fecha: \rule{2.2cm}{0.4pt}
\vspace{8pt}

\begin{tcolorbox}[colback=soft,colframe=soft,arc=2mm,left=3mm,right=3mm,top=1.5mm,bottom=1.5mm]
\small Ejercicios graduados por niveles de dificultad:\quad \estrelles{1}\ básico \qquad \estrelles{2}\ medio \qquad \estrelles{3}\ avanzado.\\[2pt] Resuélvelos en el cuaderno mostrando todos los pasos. Las rectas numéricas se pueden completar directamente en la ficha. Haz los cálculos sin calculadora.
\end{tcolorbox}

% =====================================================
\seccio{1}{Los números enteros}

\ex{1} Representa las situaciones siguientes con un número entero:\begin{tasks}(1)
  \task Un avión sobrevuela la ciudad a 432 pies de altura. \dotfill
  \task Esta noche el termómetro ha marcado 6 °C bajo cero. \dotfill
  \task El coche lo tenemos aparcado en el tercer sótano. \dotfill
  \task Sócrates nació en el año 470 antes de Cristo. \dotfill
  \task Me he gastado todo el dinero (ni tengo ni debo). \dotfill
  \task Valencia está situada a 15 m sobre el nivel del mar. \dotfill
  \task En un examen tipo test, cada respuesta incorrecta resta 1 punto. \dotfill
\end{tasks}\recorda{Primero busca el \textbf{cero} (la referencia): 0\,°C, la planta baja, el nivel del mar, no tener ni deber dinero\ldots{} Lo que está \emph{por encima} es \textbf{positivo} y lo que está \emph{por debajo}, \textbf{negativo}.}

\ex{1} Describe tres situaciones de la vida real en las que se utilizan los números negativos.

\ex{1} Escribe una situación que se pueda representar con cada entero:\begin{tasks}(5)
  \task $-12$
  \task $+8$
  \task $-1500$
  \task $+2026$
  \task $-4$
\end{tasks}

\ex{1} Clasifica estos enteros en \emph{positivos}, \emph{negativos} y \emph{ni positivos ni negativos}. ¿Cuáles son también números naturales? \[-7\qquad 12\qquad 0\qquad +3\qquad {-25}\qquad 1\qquad {-1}\qquad 100\]

\ex{2} Escribe el entero anterior y el siguiente de cada número:\begin{tasks}(5)
  \task $5$
  \task $-1$
  \task $0$
  \task $-10$
  \task $-99$
\end{tasks}

\ex{2} Escribe todos los enteros comprendidos entre:\begin{tasks}(3)
  \task $-4$ y $3$
  \task $-10$ y $-6$
  \task $-2$ y $2$
\end{tasks}

\ex{3} ¿Cuántos enteros hay comprendidos entre $-50$ y $50$? ¿Y entre $-50$ y $0$?


% =====================================================
\seccio{2}{Recta numérica, orden y comparación}

\ex{1} Representa en la recta numérica los números: \ $-2,\ 0,\ 4,\ {-1},\ 8,\ {-3},\ {-7}$ \ y \ $1$.\par\smallskip{\centering\rectabuida{-9}{9}\par}\recorda{Los \textbf{positivos} van a la \textbf{derecha} del cero y los \textbf{negativos}, a la \textbf{izquierda}. Cuanto más a la derecha, mayor es el número.}

\ex{1} Representa en la recta numérica estos números y ordénalos de menor a mayor: \[-7\qquad 3\qquad 1\qquad {-4}\qquad 6\qquad {-5}\qquad {-2}\qquad 0\]{\centering\rectabuida{-9}{9}\par}

\ex{1} Ordena los números siguientes de mayor a menor: \ $-7,\ 0,\ 1,\ {-1},\ 8,\ {-4},\ {-9},\ 3$.

\ex{1} Compara las parejas de números con el signo $<$, $>$ o $=$:\begin{tasks}(4)
  \task $-6\ \square\ 6$
  \task $0\ \square\ {-1}$
  \task $10\ \square\ {-1}$
  \task $8\ \square\ 3$
  \task $-8\ \square\ {-9}$
  \task $1\ \square\ {-2}$
  \task $-15\ \square\ {-51}$
  \task $-100\ \square\ 1$
\end{tasks}

\ex{2} Escribe un entero comprendido entre:\begin{tasks}(4)
  \task $-5$ y $-2$
  \task $-1$ y $1$
  \task $-13$ y $-11$
  \task $-100$ y $-98$
\end{tasks}

\ex{2} ¿Qué entero está\ldots\begin{tasks}(2)
  \task 5 unidades a la izquierda de $2$?
  \task 7 unidades a la derecha de $-4$?
  \task 3 unidades a la izquierda de $-6$?
  \task a 4 unidades del $0$? (hay dos)
\end{tasks}

\ex{3} ¿Verdadero o falso? Justifícalo:\begin{tasks}(1)
  \task Todo número negativo es menor que cero.
  \task $-20>-2$.
  \task Entre $-3$ y $3$ hay 5 enteros.
  \task El mayor de los números negativos es $-1$.
  \task Si $a<b$, entonces $-a>-b$.
\end{tasks}


% =====================================================
\seccio{3}{Valor absoluto y opuesto}

\ex{1} Calcula el valor absoluto de los números siguientes:\begin{tasks}(6)
  \task $\lvert -5\rvert $
  \task $\lvert -1\rvert $
  \task $\lvert 0\rvert $
  \task $\lvert 8\rvert $
  \task $\lvert -23\rvert $
  \task $\lvert 7\rvert $
\end{tasks}\recorda{El \textbf{valor absoluto} de un número es la \textbf{distancia} de ese número al 0: \ $\lvert -5\rvert =\lvert 5\rvert =5$. Nunca es negativo.}

\ex{1} Calcula el opuesto de los números siguientes:\begin{tasks}(6)
  \task $\text{op}(-8)$
  \task $\text{op}(-3)$
  \task $\text{op}(0)$
  \task $\text{op}(4)$
  \task $\text{op}(-100)$
  \task $\text{op}(15)$
\end{tasks}\recorda{El \textbf{opuesto} de un número es el que está a la \textbf{misma distancia} del cero, pero al otro lado: cambiamos el signo. \ $\text{op}(-8)=8$, \ $\text{op}(4)=-4$.}

\ex{1} Dos números opuestos se encuentran a una distancia de $36$ unidades el uno del otro. ¿Qué números son?

\ex{2} Ordena de menor a mayor los números siguientes, haciendo antes los cálculos: \[3\qquad \text{op}(4)\qquad \lvert -6\rvert \qquad 1\qquad \text{op}(5)\qquad {-2}\qquad \lvert 7\rvert \qquad \lvert -2\rvert \qquad 0\qquad 5\]

\ex{2} Compara las parejas de números:\begin{tasks}(3)
  \task $\lvert -3\rvert \ \square\ 3$
  \task $\text{op}(-2)\ \square\ {-1}$
  \task $\lvert -2\rvert \ \square\ \lvert -1\rvert $
  \task $\text{op}(\text{op}(-1))\ \square\ 0$
  \task $\lvert -4\rvert \ \square\ \text{op}(2)$
  \task $-8\ \square\ \text{op}(8)$
\end{tasks}

\ex{3} Completa los huecos con el número o los números que correspondan:\begin{tasks}(3)
  \task $\text{op}(\text{op}(\square))=-7$
  \task $\lvert \square\rvert =10$
  \task $\text{op}(\lvert \square\rvert )=-2$
  \task $\lvert \text{op}(-8)\rvert =\square$
  \task $\text{op}(\lvert -4\rvert )=\square$
  \task $\lvert \square\rvert =\text{op}(5)$
\end{tasks}

\ex{3} ¿Qué enteros tienen valor absoluto menor que $4$? ¿Y cuáles tienen valor absoluto menor o igual que $2$?


% =====================================================
\seccio{4}{Suma y resta de enteros}

\ex{1} Calcula:\begin{tasks}(5)
  \task $5-7$
  \task $-8-9$
  \task $-3+7$
  \task $4-3$
  \task $-5+1$
  \task $10-14$
  \task $-7+5$
  \task $9-13$
  \task $3+10$
  \task $9-4$
  \task $-5-5$
  \task $-4+4$
  \task $-3-2$
  \task $8+12$
  \task $-5+15$
\end{tasks}\recorda{\textbf{Mismo signo}: se suman y se pone el signo común ($-3-5=-8$). \ \textbf{Signos diferentes}: se restan (grande menos pequeño) y se pone el signo del que tiene mayor valor absoluto ($-7+4=-3$).}

\ex{1} Calcula:\begin{tasks}(6)
  \task $10-3-5$
  \task $15-9-6$
  \task $5-8+4$
  \task $9-3+5$
  \task $-2+2+7$
  \task $-10+8+6$
  \task $-10-3+8$
  \task $-4-3-2$
  \task $-1-5-7$
  \task $9-2-3$
  \task $12-4-6$
  \task $3-7+4$
  \task $5-9+8$
  \task $-13+6+4$
  \task $-2+10-15$
  \task $-11-4+8$
  \task $-5-3-4$
  \task $-8+5+6$
\end{tasks}

\ex{2} Calcula:\begin{tasks}(4)
  \task $5+7-2-4$
  \task $2-6+4-9$
  \task $9-6-7+2$
  \task $-4-5+3+8$
  \task $-8+2-7+6$
  \task $-1+5+6-7$
  \task $4-1-2-11$
  \task $-2+6-8-1$
  \task $-1+2+17-12$
  \task $-24+15+13-8$
  \task $-18-10+12+1$
  \task $11-15-13-2$
\end{tasks}\recorda{Puedes agrupar: suma todos los positivos, suma todos los negativos y resta: \ $5-9+8-2=(5+8)-(9+2)=13-11=2$.}

\ex{2} Calcula de dos maneras, de izquierda a derecha y agrupando positivos y negativos:\begin{tasks}(2)
  \task $12-4-6+9-15$
  \task $-7+3-8+10-2+5$
\end{tasks}

\ex{3} Halla el número que falta:\begin{tasks}(3)
  \task $-3+\square=5$
  \task $\square-7=-2$
  \task $4+\square=-6$
  \task $\square+9=0$
  \task $-8-\square=-3$
  \task $\square-(-4)=1$
\end{tasks}


% =====================================================
\seccio{5}{Paréntesis y corchetes}

\ex{1} Calcula, eliminando previamente los paréntesis:\begin{tasks}(3)
  \task $(-6)+(-3)$
  \task $(+7)+(-4)$
  \task $(-8)+(+1)$
  \task $(-3)+(-7)$
  \task $(+12)+(+5)$
  \task $(+3)+(-5)$
  \task $(+3)+(+4)$
  \task $(+4)+(-6)$
  \task $(+7)+(-8)$
  \task $(-5)+(+4)$
  \task $(-9)+(+3)$
  \task $(+1)+(-1)$
  \task $(-7)+(+7)$
  \task $(-12)+(-3)$
  \task $(+7)+(+5)$
  \task $(-4)+(-10)$
  \task $(+6)+(-12)$
  \task $(-9)+(+15)$
\end{tasks}\recorda{Regla de los signos para quitar paréntesis: \ $+(+\ \to\ +$ \qquad $+(-\ \to\ -$ \qquad $-(+\ \to\ -$ \qquad $-(-\ \to\ +$}

\ex{2} Calcula, eliminando previamente los paréntesis:\begin{tasks}(2)
  \task $3-(-1)$
  \task $-8+(-17)$
  \task $-4-(-7)$
  \task $-(-5)+3$
  \task $8+(-14)$
  \task $(-2)-(+3)+(-6)$
  \task $-(-5)+(-4)-(-10)$
  \task $-15-(-3)+7$
  \task $(-8)+(-3)-(-17)+(+1)$
  \task $-2-(-3)+7-8$
  \task $-(-3)+(-1)-(+3)-(-8)$
  \task $-3+(-6)-4+1$
  \task $12-25+3-8-(-6)$
  \task $-7+4-(-8)+(-3)$
  \task $-(+4)+(-8)-(-6)+4-1$
  \task $-8+(-3)+4-7-5$
  \task $-(-6)+(-1)-7+(-3)$
  \task $15-(+3)-6+3-12$
\end{tasks}

\ex{2} Calcula, resolviendo previamente los paréntesis:\begin{tasks}(4)
  \task $12+(+3-5)$
  \task $14+(+12-10)$
  \task $6-(9-7)$
  \task $15-(2-9)$
  \task $11-(-6+3)$
  \task $10-(-7-5)$
  \task $13+(-8+2)$
  \task $17+(-5-9)$
  \task $1+(-3+8)$
  \task $9-(2-1)$
  \task $-6-(2-7)$
  \task $11+(-1-3)$
  \task $6+(6+3)$
  \task $-5-(-2-8)$
  \task $3+(-8+7)$
  \task $15+(-6-8)$
\end{tasks}\recorda{Resuelve primero la operación de dentro del paréntesis, o bien quítalo: si delante hay un $-$, cambia el signo de \textbf{todos} los términos de dentro: \ $6-(9-7)=6-9+7=4$.}

\ex{2} Calcula, resolviendo previamente los paréntesis:\begin{tasks}(3)
  \task $(7-4)+(9-5)$
  \task $(2+6)+(5-8)$
  \task $(5-9)+(2-12)$
  \task $(7+3)-(5+4)$
  \task $(8-12)-(2-5)$
  \task $(10-7)-(-2-6)$
  \task $-(8+4)+(5-9)$
  \task $-(6-2)-(7-9)$
  \task $(-8-2)+(-1-5)$
  \task $(7-1)-(-5+11)$
  \task $(15-8)+(-6-12)$
  \task $(-2+3)-(8-4)$
  \task $(17-12)-(-7-8)$
  \task $(-10+3)-(2-9)$
  \task $-(6-4)+(8+1)$
  \task $-(-7-1)-(2-10)$
\end{tasks}

\ex{3} Calcula, empezando por los paréntesis de dentro:\begin{tasks}(3)
  \task $6+[5+(7+2)]$
  \task $8+[4-(3+5)]$
  \task $10-[6+(2+7)]$
  \task $15-[2-(6-10)]$
  \task $15-[10-(8+4)]$
  \task $12-[7-(2-10)]$
  \task $(-6)+[5+(2-12)]$
  \task $(-7)-[3-(4-9)]$
  \task $(2-10)+[5-(8+2)]$
  \task $(12-3)-[1-(2-6)]$
  \task $[9-(+5)]+[7+(-10)]$
  \task $[10-(-2)]-[5-(+12)]$
  \task $[8-(6+4)]-(5-7)$
  \task $[1+(6-9)]-(8-12)$
  \task $2-[6-(-1-21)]$
  \task $-1+[-5-(2-13)]$
\end{tasks}


% =====================================================
\seccio{6}{Multiplicación y división}

\ex{1} Resuelve las multiplicaciones siguientes:\begin{tasks}(3)
  \task $(-4)\cdot (+7)$
  \task $(+2)\cdot (+6)$
  \task $(-5)\cdot (-7)$
  \task $(+3)\cdot (-8)$
  \task $(-9)\cdot (-3)$
  \task $(-6)\cdot (+4)$
  \task $3\cdot (-2)$
  \task $4\cdot (+5)$
  \task $8\cdot (-6)$
  \task $-5\cdot (+3)$
  \task $-2\cdot (-4)$
  \task $-6\cdot (+3)$
  \task $(+3)\cdot (-5)\cdot (+2)$
  \task $(-4)\cdot (-1)\cdot (+6)$
  \task $(-2)\cdot (-7)\cdot (-2)$
  \task $(+5)\cdot (-4)\cdot (-3)$
  \task $(-8)\cdot 4\cdot (-6)$
  \task $5\cdot 3\cdot (-9)\cdot 4$
\end{tasks}\recorda{Regla de los signos: \ $+\cdot+\ \to\ +$ \qquad $+\cdot-\ \to\ -$ \qquad $-\cdot+\ \to\ -$ \qquad $-\cdot-\ \to\ +$. \ Para la división es igual.}

\ex{1} Calcula:\begin{tasks}(3)
  \task $(-8):(+2)$
  \task $(+20):(-10)$
  \task $(-12):(-4)$
  \task $(-4):(+2)$
  \task $(+21):(-7)$
  \task $(-12):(+6)$
  \task $(-15):(-3)$
  \task $(+32):(+8)$
  \task $(-36):(+9)$
  \task $(+42):(-7)$
  \task $(-48):(-8)$
  \task $(+54):(+6)$
  \task $[(+80):(-8)]:(-5)$
  \task $[(-70):(-2)]:(-7)$
  \task $(+50):[(-30):(+6)]$
  \task $(-40):[(+24):(+3)]$
\end{tasks}

\ex{2} Halla el número que falta:\begin{tasks}(3)
  \task $(-6)\cdot(\ \square\ )=-18$
  \task $(\ \square\ )\cdot(-3)=-24$
  \task $(\ \square\ )\cdot(-5)=+35$
  \task $(+15)\cdot(\ \square\ )=+60$
  \task $(-8)\cdot(\ \square\ )=40$
  \task $(\ \square\ )\cdot7=-77$
\end{tasks}

\ex{2} Calcula (las multiplicaciones y divisiones seguidas se hacen de izquierda a derecha):\begin{tasks}(4)
  \task $[(+6)\cdot (-4)]:(-3)$
  \task $[(-15)\cdot (-2)]:(+6)$
  \task $(-5)\cdot [(+12):(-3)]$
  \task $[(-5)\cdot (+12)]:(-3)$
  \task $24:(-2)\cdot (-3)$
  \task $(-40)\cdot [(-10):(-5)]$
  \task $(-70)\cdot (-5):(+2)$
  \task $(+100):(+4)\cdot (-2)$
\end{tasks}

\ex{3} Halla dos números enteros que cumplan, en cada caso:\begin{tasks}(1)
  \task El producto es $-12$ y la suma es $1$.
  \task El producto es $18$ y la suma es $-9$.
  \task El producto es $-20$ y la suma es $-1$.
\end{tasks}

\ex{3} Sin hacer los cálculos, di qué signo tendrá el resultado. Después, compruébalo:\begin{tasks}(1)
  \task $(-1)\cdot (-2)\cdot (-3)\cdot (-4)$
  \task $(-5)\cdot (+2)\cdot (-1)\cdot (-3)\cdot (+4)$
  \task El producto de siete factores iguales a $-1$.
\end{tasks}


% =====================================================
\seccio{7}{Operaciones combinadas}

\ex{1} Calcula:\begin{tasks}(2)
  \task $(-8)\cdot (+2)+(-5)\cdot (-3)$
  \task $(+40):(-8)-(-30):(+6)$
  \task $(-2)\cdot (-9)+(-24):(-3)-(-6)\cdot 4$
  \task $(+27):(-3)-(+3)\cdot (-5)-(-6)\cdot (-2)$
  \task $(-3)\cdot [(-2)+(-4)]$
  \task $(+4)\cdot [(-5)+(+2)]$
  \task $(+6):[(+5)-(+7)]$
  \task $(-20):[(-6)-(-2)]$
\end{tasks}\recorda{Orden de las operaciones: \ 1) paréntesis y corchetes (de dentro hacia fuera); \ 2) multiplicaciones y divisiones, de izquierda a derecha; \ 3) sumas y restas.}

\ex{2} Calcula:\begin{tasks}(2)
  \task $[(-8)+(+7)]\cdot (-3)$
  \task $[(-9)+(-3)]:(+6)$
  \task $19-(-3)\cdot [5-(+8)]$
  \task $12+(-5)\cdot [8+(-9)]$
  \task $12-[13-(-7)]:(-5)$
  \task $10-(+20):[7+(-3)]$
  \task $(-2)\cdot (5-7)-(-3)\cdot (8-6)$
  \task $(9-6)\cdot (-2)+(13+3):(-4)$
\end{tasks}

\ex{3} Calcula:\begin{tasks}(2)
  \task $-2\cdot [3-(-4+6)\cdot (-2)]+15:(-5)$
  \task $[(-3)\cdot 4-(-2)\cdot 5]:(-2)$
  \task $18:[-7-(-4)]-(-2)\cdot [5-3\cdot (-1)]$
  \task $(-4)\cdot [6-2\cdot (3-7)]-[-10:(-5)]$
\end{tasks}

\ex{3} Pon paréntesis donde haga falta para que las igualdades sean ciertas:\begin{tasks}(3)
  \task $6-4\cdot2-5=-6$
  \task $-3+5\cdot2-4=0$
  \task $10-2\cdot3+1=2$
\end{tasks}


% =====================================================
\seccio{8}{Problemas con enteros}

\ex{1} ¿Qué diferencia de temperatura soporta una persona que pasa de la cámara de conservación de las verduras, que se encuentra a $4$\,°C, a la del pescado congelado, que está a $-18$\,°C? ¿Y si pasa de la cámara del pescado a la de las verduras?\recorda{Escribe cada dato como un entero \textbf{con su signo} (\emph{sube, gana, sobre cero} $\to +$; \ \emph{baja, pierde, debe, bajo cero} $\to -$). Para calcular una \textbf{diferencia} o variación: \ \textbf{valor final} $-$ \textbf{valor inicial}.}

\ex{1} Mónica sale en el ascensor desde la planta cero de su edificio. El ascensor sube 5 plantas, después baja 3, sube 5, baja 8, sube 10, sube 5 y baja 6. ¿En qué planta está?

\ex{1} Juan le debe 14 euros a su hermana. Si le paga 9 euros, ¿cuánto le debe?

\ex{1} En una estación de esquí, el termómetro marcaba $14$\,° bajo cero a las 8 de la mañana; al mediodía la temperatura había subido 10 grados y a las 19.00 había bajado 5 grados respecto al mediodía. ¿Qué temperatura hacía a las 19.00?

\ex{1} Si estaba buceando a una profundidad de 2 metros y ahora estoy sobre un acantilado a 8 m de altura, ¿qué distancia hay entre los dos lugares?

\ex{1} El primer mes del año la gasolina bajó 5 céntimos, el segundo subió 3 céntimos, el tercero subió 8 céntimos y el cuarto bajó 2 céntimos. Exprésalo con números enteros y calcula la variación total del precio en este tiempo.

\ex{2} Un día de invierno, la temperatura a las seis de la mañana era de $3$\,° bajo cero. Al mediodía había subido 9 grados, pero a las doce de la noche el termómetro marcaba $1$\,° bajo cero. ¿Qué diferencia de temperatura ha habido entre el mediodía y la medianoche?

\ex{2} Euclides fue un matemático que vivió 60 años y murió en el año 265 a.\,C. ¿En qué año nació?

\ex{2} El Mulhacén, el pico más alto de la península ibérica, mide $3478$ metros de altura. El Sistema del Trave, una de las simas más profundas del mundo, llega a $1441$ metros de profundidad. Halla la diferencia de altitud entre los dos puntos.

\ex{2} Pablo tenía 850 euros en su cartilla de ahorros. Ha ingresado 250 euros cada mes durante los 5 últimos meses y después sacó 2300 para pagar al carpintero. ¿Qué saldo le queda?

\ex{3} Arquímedes nació en el año 287 a.\,C. y murió en el año 212 a.\,C. ¿Cuántos años vivió?

\ex{3} Un submarino se encuentra a 120 m de profundidad. Sube 45 m y después baja el triple de lo que ha subido. ¿A qué profundidad se encuentra ahora?

\ex{3} En un examen tipo test, cada respuesta correcta suma 3 puntos, cada respuesta incorrecta resta 2 y cada pregunta en blanco vale 0. Lucía ha acertado 11, ha fallado 6 y ha dejado 3 en blanco. ¿Qué puntuación ha obtenido? ¿Cuántas debería haber acertado, como mínimo, para no tener una puntuación negativa si hubiera fallado 9 y hubiera dejado 3 en blanco?

\vfill
\begin{center}\color{muted}\small\faCheckCircle[regular]\ \ ¡Ánimo! Si te cuesta algún ejercicio, repasa los apuntes interactivos de la unidad.\end{center}

\end{document}
